\section*{Capítulo \ref{ch:g12:age-of-the-earth} --- A datação radiométrica e a idade da Terra}

\begin{solution}{exo:g12:age-of-the-earth:1}
(a) um \omterm{def:g12:age-of-the-earth:closed-system}{sistema fechado}; (b) uma \omterm{def:g12:age-of-the-earth:isochron}{isócrona}; (c) a \omterm{def:g12:age-of-the-earth:age}{idade da Terra}.
\end{solution}

\begin{solution}{exo:g12:age-of-the-earth:2}
Plagioclásio (0.35), hornblenda (1.20), rocha total (3.10), feldspato
potássico (6.40), biotita (24.0). A biotita fica mais alta: ela contém mais
rubídio-87 para o seu estrôncio-86, então o seu rubídio produziu mais
estrôncio-87 desde que o granito se formou.
\end{solution}

\begin{solution}{exo:g12:age-of-the-earth:3}
$T = \ln 2/\lambda = 0.6931/\num{1.3972e-11} \approx \num{4.96e10}$ anos,
49.6 bilhões de anos: $49.6/4.54 \approx 11$ vezes a \omterm{def:g12:age-of-the-earth:age}{idade da Terra}.
\end{solution}

\begin{solution}{exo:g12:age-of-the-earth:4}
$\lambda t = \num{1.3972e-11} \times \num{3.10e8} \approx 0.004331$, e
$e^{0.004331} - 1 \approx 0.004341$. Para 480 milhões de anos, $\lambda t
\approx 0.006707$ e $e^{\lambda t} - 1 \approx 0.006729$. Os dois ficam perto
do próprio $\lambda t$, já que $\lambda t$ é pequeno.
\end{solution}

\begin{solution}{exo:g12:age-of-the-earth:5}
Phillips (1860), Kelvin (1862), Joly (1899), Holmes (1927), Patterson
(1956). $4540/100 \approx 45$ vezes.
\end{solution}

\begin{solution}{exo:g12:age-of-the-earth:6}
$0.7080 + 10.0 \times 0.004341 \approx 0.7514$ hoje; 0.7080 quando o granito
se formou, como todo mineral do granito.
\end{solution}

\begin{solution}{exo:g12:age-of-the-earth:7}
(1) $m = (0.77525 - 0.71806)/(9.40 - 0.90) = 0.05719/8.50 \approx 0.006728$;
$t = \ln(1.006728)/\num{1.3972e-11} \approx \num{4.80e8}$ anos, 480 milhões
de anos. (2) $0.71806 + 1.70 \times 0.006728 \approx 0.72950$ e $0.71806 + 4.20
\times 0.006728 \approx 0.74632$: as duas ficam sobre a reta. (3) Sim: o
granito corta o xisto, logo é mais jovem, e 310 é menos que 480 milhões de anos.
\end{solution}

\begin{solution}{exo:g12:age-of-the-earth:8}
Todos os seus pontos cairiam no mesmo lugar: um ponto não define nenhuma reta
nem nenhuma inclinação. A biotita fica muito à direita das outras; um pequeno
erro num ponto muda então a inclinação da reta que passa por ela e pelo
plagioclásio muito menos do que mudaria para dois pontos próximos.
\end{solution}

\begin{solution}{exo:g12:age-of-the-earth:9}
(1) O seu rubídio não muda, então ela conserva o seu $x$; ela tem menos
estrôncio-87, então o seu $y$ cai: o ponto desce na vertical. (2) A reta que
passa por ela e pelo plagioclásio é menos inclinada, e dá uma idade jovem
demais. (3) A biotita não fica mais sobre a reta que passa pelos outros
minerais: os pontos não estão alinhados, sinal de que uma amostra não foi um
\omterm{def:g12:age-of-the-earth:closed-system}{sistema fechado}.
\end{solution}

\begin{solution}{exo:g12:age-of-the-earth:10}
(1) $4540/96 \approx 47$ vezes. (2) Se parte do sal sai do mar, por exemplo
enterrada nos sedimentos, o mar contém menos sal do que todo o que os rios
trouxeram: dividir o sal presente pelo aporte anual dá um tempo curto
demais.
\end{solution}

\begin{solution}{exo:g12:age-of-the-earth:11}
$m = (34.86 - 10.34)/(50.28 - 9.46) = 24.52/40.82 \approx 0.601$, entre
as inclinações para 4.55 (0.590) e 4.60 bilhões de anos (0.610): cerca de
4.58 bilhões de anos.
\end{solution}

\begin{solution}{exo:g12:age-of-the-earth:12}
(1) O basalto, com a razão mais baixa, vem do manto; o granito, com a mais
alta, vem pelo menos em parte da crosta fundida. (2) O rubídio vai para o
líquido quando o manto se funde, então ele se concentrou na crosta, que
produziu mais estrôncio-87 ao longo das eras. (3) Não: minerais que se
cristalizam a partir de um mesmo magma compartilham a sua razão inicial, e
essas duas rochas não.
\end{solution}

\begin{solution}{exo:g12:age-of-the-earth:13}
(1) Rubídio-87: $2^{-4.54/49.61} \approx 0.94$, cerca de \qty{94}{\percent}
restante; urânio-238: $2^{-4.54/4.47} \approx 0.49$, cerca da metade. (2) $m =
e^{\lambda t} - 1 \approx \num{1.3972e-11} \times 10^6 \approx
\num{1.4e-5}$: mesmo a biotita do granito de Varden ($x = 24$) subiria
só $0.0003$ acima da razão inicial. Tão pouco estrôncio-87 se formou que mal
se distingue a inclinação de zero: um relógio lento convém às rochas antigas.
\end{solution}

\begin{solution}{exo:g12:age-of-the-earth:14}
(1) $4.438 + 0.6026 \times 19.0 \approx 15.89$; $15.5$ é menor em $0.39$,
cerca de \qty{2.4}{\percent}: o sedimento fica perto da reta. (2) 0.6026
fica entre 0.590 (4.55) e 0.610 (4.60): cerca de 4.58 bilhões de anos, dentro
da faixa de Patterson, de 4.48 a 4.62 bilhões de anos.
\end{solution}

\begin{solution}{exo:g12:age-of-the-earth:15}
(1) A reserva perde então só $46.7 - 24.4 = \qty{22.3}{TW}$ líquidos, então a
mesma reserva dura $46.7/22.3 \approx 2.1$ vezes mais. (2) $96 \times 2.1 \approx
200$ milhões de anos, ainda mais de vinte vezes curto demais. (3) A \omterm{def:g12:internal-heat:convection}{convecção
do manto}, que traz calor das profundezas e mantém o gradiente perto da
superfície alto por muito mais tempo que a condução sozinha.
\end{solution}

\begin{solution}{pb:g12:age-of-the-earth:1}
\textbf{1.} A biotita ($x = 24.0$).

\textbf{2.} Eles se cristalizaram a partir de um mesmo magma, e os minerais
que se cristalizam a partir de um mesmo líquido incorporam o seu estrôncio
com a sua razão de isótopos.

\textbf{3.} O estrôncio-86 é estável e não é produzido por nenhum
decaimento: num \omterm{def:g12:age-of-the-earth:closed-system}{sistema fechado}, o seu número não muda, então ele é uma
referência fixa.

\textbf{4.} O estrôncio-87, produzido pelo decaimento do rubídio-87; o
rubídio-87 diminuiu e o estrôncio-86 continuou o mesmo.

\textbf{5.} Uma reta horizontal em $y = 0.7080$: mesma razão, $x$ diferentes.

\textbf{6.} $(0.81218 - 0.70952)/(24.0 - 0.35) = 0.10266/23.65 \approx
0.004341$.

\textbf{7.} $(0.72146 - 0.70952)/(3.10 - 0.35) = 0.01194/2.75 \approx
0.004342$.

\textbf{8.} Elas concordam com uma diferença de 1 em 4000. A primeira é a
mais confiável: os seus dois pontos estão afastados, então o arredondamento
das razões pesa menos.

\textbf{9.} $0.70952 - 0.35 \times 0.004341 \approx 0.7080$.

\textbf{10.} $0.7080 + 6.40 \times 0.004341 \approx 0.73578$: exatamente o
valor medido.

\textbf{11.} $e^{\lambda t} = 1 + m = 1.004341$.

\textbf{12.} $\lambda t = \ln 1.004341 \approx 0.004331$.

\textbf{13.} $t = 0.004331/\num{1.3972e-11} \approx \num{3.10e8}$ anos.

\textbf{14.} Sim: o granito corta o xisto, então, pelo princípio das
relações de corte, ele é mais jovem, e 310 é menos que 480 milhões de
anos.

\textbf{15.} 310 milhões de anos ficam entre 358.9 e 298.9 milhões de anos:
o Carbonífero.

\textbf{16.} A razão $^{87}\mathrm{Sr}/^{86}\mathrm{Sr}$ do magma quando
o granito se cristalizou.

\textbf{17.} O granito: $0.7080 > 0.7035$.

\textbf{18.} O rubídio vai para o líquido quando o manto se funde, então ele
se concentrou na crosta; com mais rubídio, a crosta produziu mais
estrôncio-87.

\textbf{19.} Que o magma veio, pelo menos em parte, da \omterm{def:g12:continental-crust:continental-crust}{crosta continental}
fundida, e não só do manto.

\textbf{20.} O granito de Varden tem \textbf{310 milhões de anos}.
\end{solution}
